\section{Electrochemical--Thermal--Aging Model}
\label{sec:model_section}

\subsection{Electrochemical Model}
\label{sec:model}

The single particle model with electrolyte dynamics (SPMe) describes the
battery's electrochemical response~\cite{refSPMe}. It accounts for lithium
diffusion in representative solid particles of both electrodes and transport
through the electrolyte. This reduced model retains the electrochemical
dynamics needed for charging simulation at a lower computational cost than a
full porous-electrode model.

The terminal voltage is calculated as~\cite{refSPMe}
\begin{equation}
U=U_{\mathrm{eq}}+\eta_r+\eta_d+\Delta\Phi
\label{eq:voltage}
\end{equation}
where $U_{\mathrm{eq}}$ is the equilibrium potential, $\eta_r$ is the ohmic
overpotential, $\eta_d$ is the concentration overpotential, and $\Delta\Phi$
is the average potential difference in the electrolyte. SOC is obtained from
the solid-phase lithium concentration. The SPMe equations and their derivation
are given in~\cite{refSPMe}.
% AUTHOR_CHECK M1: Verify the voltage decomposition and the meanings/signs of
% eta_r, eta_d, and DeltaPhi against the model actually implemented.

\subsection{Thermal Model}

A lumped thermal model gives the cell temperature:
\begin{equation}
\rho_{\mathrm{eff}}\frac{\partial T}{\partial t}=\bar Q-\frac{hA}{V}(T-T_a)
\label{eq:thermal_balance}
\end{equation}
Here, $\rho_{\mathrm{eff}}$ is the effective volumetric heat capacity,
$\bar Q$ is the average volumetric heat-generation rate, $h$ is the
heat-transfer coefficient, $A$ is the heat-dissipation area, $V$ is the cell
volume, and $T_a$ is the ambient temperature. The effective heat capacity is
\begin{equation}
\rho_{\mathrm{eff}}=\frac{\sum_k\rho_kC_{p,k}L_k}{\sum_kL_k}
\label{eq:rho_eff}
\end{equation}
where $\rho_k$, $C_{p,k}$, and $L_k$ denote density, specific heat capacity,
and thickness. The sum includes the two current collectors, the two
electrodes, and the separator.

The thermal objective is the maximum temperature rise during charging:
\begin{equation}
T_p(\theta)=\max_{0\le t\le t_c(\theta)}[T(t)-T_0]
\label{eq:peak_temp}
\end{equation}
where $T_0$ is the initial cell temperature.

\subsection{Aging Model}
\label{sec:aging}

Charging-induced degradation is evaluated with the semi-empirical
capacity-loss model of Suri and Onori~\cite{refAgingModel}:
\begin{equation}
Q_s(\xi,Ah)=\sigma_{\mathrm{fun}}(\xi)Ah^z
\label{eq:aging}
\end{equation}
The vector $\xi=(\mathrm{SOC},I,T)$ contains the average SOC, current, and
temperature during charging; $Ah$ is the charge throughput; and $z$ is the
power-law exponent. The severity factor is
% AUTHOR_CHECK M2: The uploaded equation is preserved below. Its denominator
% uses 273.15+T, whereas the thermal settings use kelvin. Verify the conversion
% in the implementation before choosing a Celsius or kelvin notation here.
\begin{equation}
\sigma_{\mathrm{fun}}(\xi)=(\alpha\,\mathrm{SOC}+\beta)\exp\!\left(\frac{-E_a+\eta_a I}{R_{\mathrm{gas}}(273.15+T)}\right)
\label{eq:severity}
\end{equation}
where $\alpha$ and $\beta$ describe SOC dependence, $\eta_a$ describes current
dependence, $E_a$ is the activation energy, and $R_{\mathrm{gas}}$ is the
universal gas constant. We use $Q_s$ as a relative degradation indicator for
one charging process, rather than as a prediction of battery end-of-life.
% AUTHOR_CHECK M3: Supply the coefficients, averaging/aggregation rule, and
% justification for using this aging model with each battery parameterization.
